\section*{第 \ref{ch:g3:stages-of-life} 章——生命的各个阶段}

\begin{solution}{exo:g3:stages-of-life:1}
婴儿、儿童、青少年、\omterm{def:g3:stages-of-life:stages}{成体}、\omterm{def:g3:stages-of-life:stages}{老年期}。
\end{solution}

\begin{solution}{exo:g3:stages-of-life:2}
（a）大约从六岁起进入儿童期。（b）处于\omterm{def:g3:stages-of-life:stages}{成体}阶段。
（c）处于婴儿期——第一年是所有阶段中生长最快的。
\end{solution}

\begin{solution}{exo:g3:stages-of-life:3}
婴儿期和青少年期各有一次快速生长，两者之间是儿童期平稳的
缓慢生长。
\end{solution}

\begin{solution}{exo:g3:stages-of-life:4}
曲线变平：年龄仍沿底部增长，但
身高不再上升。
\end{solution}

\begin{solution}{exo:g3:stages-of-life:5}
在 $6$ 岁至 $8$ 岁之间：身高由约 $\qty{115}{cm}$ 增至约
$\qty{127}{cm}$——两年共增长约 $\qty{12}{cm}$。
仅第一年：身高由约 $\qty{50}{cm}$ 增至约
$\qty{75}{cm}$——约增长 $\qty{25}{cm}$，
是这两个儿童期年份合计增长量的两倍。
\end{solution}

\begin{solution}{exo:g3:stages-of-life:6}
因为只有测量方法相同，比较才公平。
第一天穿鞋、第二天穿袜子，或者书本歪斜、随意摆放，
都会造成（或掩盖）一种并未发生的“生长”。
这样的比较就失去了意义。
\end{solution}

\begin{solution}{exo:g3:stages-of-life:7}
橡子（\omterm{def:g2:from-seed-to-plant:seed}{种子}）、\omterm{prop:g2:from-seed-to-plant:cycle}{幼苗}、幼体橡树、开花结实的\omterm{def:g3:stages-of-life:stages}{成体}
橡树、老年橡树——以及，经过一段时间，死亡。
\end{solution}

\begin{solution}{exo:g3:stages-of-life:8}
开放性答案。九岁的孩子处于儿童阶段；
青少年期的快速生长尚未到来。
\end{solution}

\begin{solution}{exo:g3:stages-of-life:9}
苍蝇（数周）、猫（约一年）、人（近二十年），
橡树（数十年）。四者的阶段顺序相同，
不能跳过，也不能倒退：幼体、\omterm{def:g3:stages-of-life:stages}{成体}、
\omterm{def:g3:stages-of-life:stages}{老年期}。
\end{solution}

\begin{solution}{exo:g3:stages-of-life:10}
不对。他九岁时处于儿童期平稳的阶段——
每\emph{年}约增长 $\qty{6}{cm}$，即每月约半厘米：
相隔一个月的两次测量之间难以察觉。
曲线稳步上升，只是速度较慢；他只需
在相隔数月的两个数据点之间等待即可。
\end{solution}
